\documentclass[12pt,a4paper]{article}
\usepackage[margin=1in]{geometry}
\usepackage{booktabs}
\usepackage{longtable}
\usepackage{graphicx}
\usepackage{float}
\usepackage{amsmath}
\usepackage{hyperref}
\usepackage{xcolor}
\usepackage{caption}
\usepackage{array}
\usepackage{multirow}

\title{\textbf{Lab 3: Cache Simulator --- Analysis Report}}
\author{Sharique}
\date{\today}

\begin{document}
\maketitle
\tableofcontents
\newpage

%=============================================================================
\section{Task 1: Trace File Generation}
%=============================================================================

Ten trace files were generated, each with 500{,}000 memory accesses, using the \texttt{trace\_generator} tool. The default write ratio is 30\% (70\% reads, 30\% writes) unless overridden.

\begin{table}[H]
\centering
\caption{Generated Trace Files}
\begin{tabular}{clll}
\toprule
\textbf{\#} & \textbf{Trace Name} & \textbf{Pattern} & \textbf{Key Parameter} \\
\midrule
1 & sequential\_500k & Sequential & stride=64 \\
2 & random\_500k & Random & --- \\
3 & strided64\_500k & Strided & stride=128 \\
4 & looping\_2x4\_500k & Looping & 2 loops, 4 reps \\
5 & mixed\_high\_loc\_500k & Mixed & locality=0.8 \\
6 & mixed\_low\_loc\_500k & Mixed & locality=0.2 \\
7 & write\_intensive\_500k & Mixed & write ratio=0.8 \\
8 & low\_write\_500k & Mixed & write ratio=0.2 \\
9 & seq\_4proc\_500k & Sequential & 4 processors \\
10 & rand\_4proc\_500k & Random & 4 processors \\
\bottomrule
\end{tabular}
\end{table}

%=============================================================================
\section{Task 2: Baseline L1 Cache Analysis}
%=============================================================================

\textbf{Configuration:} L1 = 32KB, 1-way (direct-mapped), 64-byte blocks, LRU, Write-Back.

\begin{table}[H]
\centering
\caption{Task 2 --- Baseline Results (1-way, 64B blocks, LRU, Write-Back)}
\resizebox{\textwidth}{!}{%
\begin{tabular}{lrrrrrr|rrrr}
\toprule
\textbf{Trace} & \textbf{Reads} & \textbf{Writes} & \textbf{Hits} & \textbf{Misses} & \textbf{WB} & \textbf{Hit\%} & \textbf{Comp.} & \textbf{Cap.} & \textbf{Conf.} & \textbf{Coh.} \\
\midrule
Sequential     & 349,820 & 150,180 & 0       & 500,000 & 150,034 & 0.00  & 512 & 0 & 499,488 & 0 \\
Random         & 349,240 & 150,760 & 15,715  & 484,285 & 149,119 & 3.14  & 512 & 0 & 483,773 & 0 \\
Strided-64     & 350,614 & 149,386 & 0       & 500,000 & 149,306 & 0.00  & 256 & 0 & 499,744 & 0 \\
Looping        & 350,178 & 149,822 & 499,998 & 2       & 0       & 100.00& 2   & 0 & 0       & 0 \\
Mixed Hi-Loc   & 349,812 & 150,188 & 121,793 & 378,207 & 117,146 & 24.36 & 512 & 0 & 377,695 & 0 \\
Mixed Lo-Loc   & 349,466 & 150,534 & 120,988 & 379,012 & 117,655 & 24.20 & 512 & 0 & 378,500 & 0 \\
Write-Heavy    & 100,136 & 399,864 & 119,598 & 380,402 & 305,574 & 23.92 & 512 & 0 & 379,890 & 0 \\
Read-Heavy     & 400,142 & 99,858  & 120,738 & 379,262 & 79,880  & 24.15 & 512 & 0 & 378,750 & 0 \\
Seq 4-Proc     & 349,660 & 150,340 & 0       & 500,000 & 150,179 & 0.00  & 512 & 0 & 499,488 & 0 \\
Rand 4-Proc    & 350,505 & 149,495 & 15,659  & 484,341 & 147,898 & 3.13  & 512 & 0 & 483,829 & 0 \\
\bottomrule
\end{tabular}%
}
\end{table}

\textbf{Observation:} The direct-mapped cache performs terribly for most workloads. Sequential and strided patterns yield 0\% hit rate because every access maps to a new block and the working set far exceeds the 32KB cache. The looping pattern achieves a perfect 100\% hit rate since only 2 unique addresses fit trivially in cache. Mixed patterns hover around 24\% due to partial spatial reuse. Over 99.9\% of all misses are \textbf{conflict misses}, highlighting the fundamental limitation of direct-mapped caches. Write ratio doesn't affect hit rate but significantly impacts write-backs (305K for 80\% writes vs.\ 80K for 20\% writes).

%=============================================================================
\section{Task 3: Varying Associativity (2, 4, 8-way)}
%=============================================================================

\textbf{Configuration:} Same as Task 2 but associativity changed to 2, 4, and 8.

\begin{table}[H]
\centering
\caption{Task 3 --- Hit Ratio (\%) by Associativity}
\begin{tabular}{lrrrr}
\toprule
\textbf{Trace} & \textbf{1-way} & \textbf{2-way} & \textbf{4-way} & \textbf{8-way} \\
\midrule
Sequential     & 0.00  & 0.00  & 0.00  & 0.00  \\
Random         & 3.14  & 3.15  & 3.15  & 3.15  \\
Strided-64     & 0.00  & 0.00  & 0.00  & 0.00  \\
Looping        & 100.00& 100.00& 100.00& 100.00\\
Mixed Hi-Loc   & 24.36 & 24.59 & 25.38 & 25.39 \\
Mixed Lo-Loc   & 24.20 & 24.42 & 25.32 & 25.32 \\
Write-Heavy    & 23.92 & 24.09 & 25.28 & 25.30 \\
Read-Heavy     & 24.15 & 24.35 & 25.31 & 25.33 \\
Seq 4-Proc     & 0.00  & 0.00  & 0.00  & 0.00  \\
Rand 4-Proc    & 3.13  & 3.14  & 3.14  & 3.12  \\
\bottomrule
\end{tabular}
\end{table}

\textbf{Observation:} Increasing associativity provides \textbf{marginal improvement} ($\sim$1\%) for mixed patterns by reducing conflict misses. The sequential and strided patterns see no benefit because their working sets far exceed the cache size---the misses are fundamentally due to the data not fitting, not due to conflicts within a set. The looping pattern remains perfect. The biggest gain is from 1-way to 4-way for mixed patterns ($\sim$24\% $\to$ 25.3\%). Beyond 4-way, diminishing returns are evident (4-way vs.\ 8-way shows $<$0.1\% difference).

%=============================================================================
\section{Task 4: Varying Block Size (16, 32, 128 bytes)}
%=============================================================================

\textbf{Configuration:} Same as Task 2 but block size changed to 16, 32, and 128.

\begin{table}[H]
\centering
\caption{Task 4 --- Hit Ratio (\%) by Block Size}
\begin{tabular}{lrrrr}
\toprule
\textbf{Trace} & \textbf{16B} & \textbf{32B} & \textbf{64B (base)} & \textbf{128B} \\
\midrule
Sequential     & 0.00  & 0.00  & 0.00  & 50.00 \\
Random         & 3.14  & 3.14  & 3.14  & 3.16  \\
Strided-64     & 0.00  & 0.00  & 0.00  & 0.00  \\
Looping        & 100.00& 100.00& 100.00& 100.00\\
Mixed Hi-Loc   & 24.36 & 24.36 & 24.36 & 27.21 \\
Mixed Lo-Loc   & 24.20 & 24.20 & 24.20 & 27.38 \\
Write-Heavy    & 23.92 & 23.92 & 23.92 & 27.11 \\
Read-Heavy     & 24.15 & 24.15 & 24.15 & 27.17 \\
Seq 4-Proc     & 0.00  & 0.00  & 0.00  & 50.00 \\
Rand 4-Proc    & 3.13  & 3.13  & 3.13  & 3.12  \\
\bottomrule
\end{tabular}
\end{table}

\textbf{Observation:} Block size has a dramatic effect on sequential access. With 128B blocks, the sequential trace jumps from 0\% to \textbf{50\%} hit rate because each 128B block now covers two consecutive 64-byte-stride accesses---the second access in every pair is a hit. Mixed patterns gain $\sim$3\% from the larger block size. However, random and strided-64 patterns see negligible change because larger blocks don't help when access patterns don't exhibit spatial locality within the block. Notably, 16B and 32B blocks give \textbf{identical results} to 64B for this particular trace set, suggesting the trace generator's stride perfectly aligned with block boundaries.

%=============================================================================
\section{Task 5: Varying Replacement Policy}
%=============================================================================

\textbf{Configuration:} 32KB, \textbf{4-way}, 64B blocks, Write-Back (4-way used instead of 1-way because replacement policy has no effect on direct-mapped caches---there's only one slot per set).

\begin{table}[H]
\centering
\caption{Task 5 --- Hit Ratio (\%) by Replacement Policy (4-way)}
\begin{tabular}{lrrrr}
\toprule
\textbf{Trace} & \textbf{LRU} & \textbf{FIFO} & \textbf{PLRU} & \textbf{Random} \\
\midrule
Sequential     & 0.00  & 0.00  & 0.00  & 0.00  \\
Random         & 3.14  & 3.13  & 3.10  & 3.14  \\
Strided-64     & 0.00  & 0.00  & 0.00  & 0.00  \\
Looping        & 100.00& 100.00& 100.00& 100.00\\
Mixed Hi-Loc   & 25.38 & 24.36 & 21.62 & 24.36 \\
Mixed Lo-Loc   & 25.30 & 24.18 & 21.24 & 24.21 \\
Write-Heavy    & 25.28 & 23.90 & 20.31 & 23.90 \\
Read-Heavy     & 25.31 & 24.12 & 21.02 & 24.13 \\
Seq 4-Proc     & 0.00  & 0.00  & 0.00  & 0.00  \\
Rand 4-Proc    & 3.14  & 3.14  & 3.12  & 3.13  \\
\bottomrule
\end{tabular}
\end{table}

\textbf{Observation:} For workloads with some locality (mixed patterns), the ranking is: \textbf{LRU > FIFO $\approx$ Random > PLRU}. LRU outperforms all others by $\sim$1--4\% because it makes the best eviction decisions based on recency. Surprisingly, PLRU performs worst ($\sim$21\%) because its tree-based approximation makes poor eviction choices under this access pattern. FIFO and Random perform similarly ($\sim$24\%). For streaming patterns (sequential, strided), all policies are equally ineffective since every access is a miss regardless of eviction strategy.

%=============================================================================
\section{Task 6: Write-Through Policy}
%=============================================================================

\textbf{Configuration:} Same as Task 2 but using Write-Through instead of Write-Back.

\begin{table}[H]
\centering
\caption{Task 6 --- Write-Through vs.\ Write-Back Comparison}
\begin{tabular}{lr|rr}
\toprule
\textbf{Trace} & \textbf{Hit\%} & \textbf{WB (Write-Back)} & \textbf{WB (Write-Through)} \\
\midrule
Sequential     & 0.00  & 150,034 & 150,076 \\
Random         & 3.14  & 149,119 & 145,930 \\
Strided-64     & 0.00  & 149,306 & 149,306 \\
Looping        & 100.00& 0       & 0       \\
Mixed Hi-Loc   & 24.36 & 117,146 & 113,635 \\
Mixed Lo-Loc   & 24.20 & 117,655 & 113,967 \\
Write-Heavy    & 23.92 & 305,574 & 303,976 \\
Read-Heavy     & 24.15 & 79,880  & 75,798  \\
\bottomrule
\end{tabular}
\end{table}

\textbf{Observation:} The hit ratio is \textbf{identical} between Write-Back and Write-Through because the write policy doesn't affect whether data is found in cache. The difference is in \textbf{memory traffic}: Write-Through sends every write directly to memory (more traffic but simpler consistency), while Write-Back defers writes until eviction (fewer memory writes but requires dirty-bit tracking). The write-back counts are slightly lower for Write-Through because dirty blocks don't accumulate---data is always clean, so evictions don't require write-backs. However, Write-Through still generates write-backs for blocks that were brought in on a read miss and later written.

%=============================================================================
\section{Task 7: Prefetching}
%=============================================================================

\textbf{Configuration:} Same as Task 2 with prefetching enabled (distance=4).

\begin{table}[H]
\centering
\caption{Task 7 --- With Prefetching vs.\ Baseline}
\begin{tabular}{lrr|rr}
\toprule
\textbf{Trace} & \textbf{Hit\% (base)} & \textbf{Hit\% (pref)} & \textbf{Misses (base)} & \textbf{Misses (pref)} \\
\midrule
Sequential     & 0.00  & 41.17 & 500,000 & 294,172 \\
Random         & 3.14  & 3.14  & 484,285 & 484,285 \\
Strided-64     & 0.00  & 41.24 & 500,000 & 293,820 \\
Looping        & 100.00& 100.00& 2       & 2       \\
Mixed Hi-Loc   & 24.36 & 24.39 & 378,207 & 378,055 \\
Mixed Lo-Loc   & 24.20 & 24.23 & 379,012 & 378,875 \\
Write-Heavy    & 23.92 & 23.93 & 380,402 & 380,351 \\
Read-Heavy     & 24.15 & 24.18 & 379,262 & 379,093 \\
Seq 4-Proc     & 0.00  & 41.14 & 500,000 & 294,278 \\
Rand 4-Proc    & 3.13  & 3.13  & 484,341 & 484,341 \\
\bottomrule
\end{tabular}
\end{table}

\textbf{Observation:} Prefetching has a \textbf{massive impact on sequential/strided patterns}, improving hit rate from 0\% to $\sim$41\% and reducing misses by $\sim$41\%. The prefetcher correctly predicts the next blocks in the stride pattern and fetches them before they're needed. For \textbf{random access}, prefetching has \textbf{zero effect}---the prefetched blocks are never used because the next access goes to an unpredictable location. Mixed patterns see negligible improvement ($<$0.1\%) because their access patterns are not consistently predictable enough for the stride prefetcher to help.

%=============================================================================
\section{Task 8: Adding L2 Cache}
%=============================================================================

\textbf{Configuration:} L1 same as Task 2. L2 = 256KB, 16-way, 64B blocks, PLRU, Write-Back.

\begin{table}[H]
\centering
\caption{Task 8 --- L1 and L2 Statistics}
\begin{tabular}{lrr|rr}
\toprule
\textbf{Trace} & \textbf{L1 Hit\%} & \textbf{L1 Misses} & \textbf{L2 Hit\%} & \textbf{L2 Miss\%} \\
\midrule
Sequential     & 0.00  & 500,000 & 43.57 & 56.43 \\
Random         & 3.14  & 484,285 & 61.92 & 38.08 \\
Looping        & 100.00& 2       & 50.00 & 50.00 \\
Mixed Hi-Loc   & 24.36 & 378,207 & 50.71 & 49.29 \\
\bottomrule
\end{tabular}
\end{table}

\textbf{Observation:} The L2 cache successfully catches a significant portion of L1 misses. For \textbf{random access}, the L2 hit rate is 61.92\%---this is because the random trace's working set (16,320 unique addresses $\times$ 64B = $\sim$1MB) partially fits in the 256KB L2 with 16-way associativity reducing conflicts. For \textbf{sequential}, the L2 captures 43.57\% of L1 misses. The L1 hit rates remain unchanged (the L2 doesn't affect L1 behavior). The key takeaway: multi-level caches reduce the \textbf{effective miss penalty} by catching misses at L2 before they go to slow main memory.

%=============================================================================
\section{Task 9: MESI Coherence Protocol Assessment}
%=============================================================================

\textbf{Configuration:} Baseline L1 with multiprocessor traces (4 processors).

\begin{table}[H]
\centering
\caption{Task 9 --- MESI State Transitions (Sequential 4-Processor Trace)}
\begin{tabular}{lrrrr}
\toprule
\textbf{From $\backslash$ To} & \textbf{Modified} & \textbf{Exclusive} & \textbf{Shared} & \textbf{Invalid} \\
\midrule
Modified  & 0      & 104,982 & 0 & 0 \\
Exclusive & 105,000& 0       & 0 & 0 \\
Shared    & 0      & 0       & 0 & 0 \\
Invalid   & 143    & 369     & 0 & 0 \\
\bottomrule
\end{tabular}
\end{table}

\begin{table}[H]
\centering
\caption{Task 9 --- MESI State Transitions (Random 4-Processor Trace)}
\begin{tabular}{lrrrr}
\toprule
\textbf{From $\backslash$ To} & \textbf{Modified} & \textbf{Exclusive} & \textbf{Shared} & \textbf{Invalid} \\
\midrule
Modified  & 0      & 103,612 & 0 & 0 \\
Exclusive & 103,638& 0       & 0 & 0 \\
Shared    & 0      & 0       & 0 & 0 \\
Invalid   & 143    & 369     & 0 & 0 \\
\bottomrule
\end{tabular}
\end{table}

\textbf{Assessment of MESI correctness:}

The observed transitions are consistent with MESI protocol behavior for a single-cache simulation:

\begin{itemize}
    \item \textbf{I $\to$ E} (369 transitions): When a processor first reads a block, it enters the \textit{Exclusive} state since no other cache holds it.
    \item \textbf{I $\to$ M} (143 transitions): When a processor first writes to a block not in any cache, it goes directly to \textit{Modified}.
    \item \textbf{E $\to$ M} ($\sim$105K transitions): When a processor writes to an Exclusive block, it silently upgrades to \textit{Modified} without bus traffic (this is the key advantage of MESI over MSI).
    \item \textbf{M $\to$ E} ($\sim$105K transitions): When a modified block is written back to memory (on eviction), it transitions back to Exclusive on the next load.
    \item \textbf{No Shared state}: Since the simulator runs a single shared cache for all processors (not per-processor caches), there are no S $\to$ I invalidations or M $\to$ S downgrades. True multiprocessor coherence traffic would appear with per-processor cache simulation.
\end{itemize}

%=============================================================================
\section{Summary of Key Findings}
%=============================================================================

\begin{enumerate}
    \item \textbf{Direct-mapped caches are extremely conflict-prone.} Over 99.9\% of misses across all non-looping patterns are conflict misses.
    \item \textbf{Associativity helps, but with diminishing returns.} The biggest gain is 1-way $\to$ 4-way ($\sim$1\% improvement for mixed patterns). Beyond 4-way, improvements are negligible.
    \item \textbf{Block size dramatically affects spatial-locality workloads.} 128B blocks improve sequential hit rate from 0\% to 50\%.
    \item \textbf{LRU is the best replacement policy}, outperforming FIFO, Random, and PLRU by 1--4\% on mixed workloads.
    \item \textbf{Write policy doesn't affect hit rate} but affects memory traffic. Write-Back is preferred for reducing bandwidth.
    \item \textbf{Prefetching is highly effective for predictable patterns} (41\% improvement for sequential) but useless for random access.
    \item \textbf{L2 cache catches 44--62\% of L1 misses}, significantly reducing the effective miss penalty.
    \item \textbf{MESI transitions are correct:} I$\to$E for first reads, E$\to$M for writes (silent upgrade), M$\to$E on write-back.
\end{enumerate}

\end{document}
